% TODO: styling
\mdtheorem[style=theoremstyle,frametitlebackgroundcolor=black,frametitlefontcolor=white]{exercise}{Exercise}

\newcommand{\llpr}[1]{\llparenthesis #1 \rrparenthesis}
\newcommand{\FamSet}{\mathsf{Fam}(\mathsf{Set})}

\chapter{Type Refinements as Indexed Inductive Types}

In the course of writing code in a simply typed language, one may find oneself wishing to constrain
the values of a type based on some ``feature'' of its values. For instance, when writing functions
over lists or trees, it may be desirable to enforce constraints on the length or height of the
input, or to ensure that the values it stores conform to some precondition. A classic example is
the list \texttt{zip} operation: the author of this function may wish to enforce that its arguments
have the same length.

\emph{Refinement types} enrich simple type systems to allow programmers to express such constraints.
For example, the type of lists of length \texttt{n} can be written
\[ \texttt{\{xs : List A | length(xs) = n\}} \]

In dependently typed languages, such refinement types can be na\"ively translated as dependent
pairs, adjoining to values of the list type proofs that they satisfy the length constraint:

\[ \sum_{xs ~:~ \texttt{List}~\alpha} \texttt{length}(xs) = n \]

However, dependent types permit a more elegant encoding: such constraints can be directly ``baked
into'' the definition of an inductive type via type indices, ensuring that all values inhabiting the
type are correct by construction. Our example above, for instance, yields the type of length-indexed
vectors:\footnote{Throughout this chapter, we will write code snippets in a syntax similar to that
of Lean 4, though we assume all constructors live in the global namespace. Those familiar with Agda
can read \texttt{inductive} as \texttt{data} and \texttt{Type} as \texttt{Set}.}

\begin{lstlisting}
inductive Vector (A : Type) : Nat -> Type
  | nil : Vector A 0
  | cons : {n : Nat} -> A -> Vector A n -> Vector A (n + 1)
\end{lstlisting}

The type \texttt{Vector} naturally internalizes the \texttt{length} refinement from the simply typed
setting. In this chapter, we will use a categorical construction to generalize this transformation
to arbitrary inductive types and a broad class of refinements. In other words, we answer the
following question: how can a refinement of a type in a simply typed language be translated into an
indexed inductive type in a dependently typed one? Along the way, we will see how to translate
familiar concepts like inductive datatypes and their folds into categorical language. The content of
this chapter is based on the work of Atkey, Johann, and Ghani \citep{atkey2011refining}.

% Note that in this chapter, we will tacitly work over a set-theoretic semantics of a basic lambda
% calculus. However, to avoid having to stray too far afield into a general account of (dependent)
% inductive types, we will opt not to give a formal grammar for our syntax, and instead rely on code
% snippets to characterize the entities of interest.

\section{Categorifying Inductive Types}

We will begin by adding a new programming-language feature to our categorical repertoire:
definitions of inductive datatypes.

\subsection{Wherefore More Constructions?}

Let us first consider why the facilities afforded to us by our existing semantic development are
inadequate to express inductive type definitions in full generality. This may not be immediately
obvious: for instance, as we have previously seen, the datatype of booleans
\begin{lstlisting}
inductive Bool : Type where
  | true : Bool
  | false : Bool
\end{lstlisting}
has elements isomorphic to those of $\mathtt{Unit} \oplus \mathtt{Unit}$, which we interpret
categorically as $1 \oplus 1$. We can even model finite parameterized datatypes: the option datatype
\begin{lstlisting}
inductive Option (A : Type) : Type where
  | some : A -> Option A
  | none : Option A
\end{lstlisting}
has elements isomorphic to those of $\texttt{A} \oplus \texttt{Unit}$, which we interpret as
$\llbr{\texttt{A}} \oplus 1$ (assuming that we are able to interpret \texttt{A}).

You may notice that our wording above is somewhat roundabout: this is intentional. An inductive type
is characterized not merely by its ground inhabitants, but also by the \emph{maps into} the type
(i.e., its constructor forms). An inductive type also comes equipped with an \emph{iterator} (or
\emph{fold}) that admits recursive definitions over the type. The account above does not
characterize these facets of an inductive definition.

Moreover, we also currently lack the facilities to provide a general categorical account of infinite
inductive types, such as the type of natural numbers
\begin{lstlisting}
inductive Nat : Type where
  | zero : Nat
  | succ : Nat -> Nat\end{lstlisting}
or the type of lists
\begin{lstlisting}
inductive List (A : Type) : Type where
  | nil  : List A
  | cons : A -> List A -> List A
\end{lstlisting}
since the tools we have previously used to give the semantics of lambda calculi---(finite)
(co)products, exponential objects, and initial and terminal objects---do not (at least
straightforwardly)\footnote{We would prefer a more ``natural'' account of inductive types than, say,
Church encodings.} admit the construction of sums and products of unbounded arity. To give a
complete account of general inductive definitions, we will need to develop some new (and repurpose
some old) tools.

We will motivate our translation of inductive datatypes using our standard approach to translating
PL concepts into category theory: translating inference rules. Specifically, we will translate the
inference rules for booleans and natural numbers into categorical constructions, then observe some
patterns that will allow us to generalize these translations to arbitrary inductive types.

\subsection{Categorifying Booleans}

To begin, we return to the example of the type \texttt{Bool}. Its constructors are characterized by
the following typing rules:

\begin{mathpar}
\inferrule*[right=T-true]{~}{\vdash \texttt{true} : \texttt{Bool}}

\inferrule*[right=T-false]{~}{\vdash \texttt{false} : \texttt{Bool}}
\end{mathpar}

Note that we write these rules with an empty context, rather than an arbitrary $\Gamma$. Recalling
that the empty context is represented by the terminal object $1$, these rules correspond to the
existence of two \emph{global elements} $\texttt{true}, \texttt{false} : 1 \to \llbr{\texttt{Bool}}$.

\begin{definition}[Global element]
  In a category with terminal objects, a \emph{global element} of an object $A$ is a morphism
  $1 \to A$. It is the categorical interpretation of an ``element'' of an object.
\end{definition}

Here, our use of global elements corresponds to the fact that \texttt{true} and \texttt{false} are
distinguished members of the ``collection'' denoted by \texttt{Bool}. The rules above translate into
the following diagram:

% https://q.uiver.app/#q=WzAsMyxbMCwwLCIxIl0sWzIsMCwiXFx0ZXh0dHR7Qm9vbH0iXSxbNCwwLCIxIl0sWzAsMSwiXFx0ZXh0dHR7dHJ1ZX0iXSxbMiwxLCJcXHRleHR0dHtmYWxzZX0iLDJdXQ==
\[\begin{tikzcd}
	1 && {\llbr{\texttt{Bool}}} && 1
	\arrow["{\texttt{true}}", from=1-1, to=1-3]
	\arrow["{\texttt{false}}"', from=1-5, to=1-3]
\end{tikzcd}\]

However, we want to be able to map \emph{out} of our inductive types as well. To do so, we must
define the \emph{iterator} for the Boolean type, canonically known as \texttt{if}:\footnote{We will
adopt the convention of writing iterators with their \emph{major premise} (i.e., the value being
``iterated'') as the final argument. This is opposite the usual notation in functional programming
but will better agree with the categorical formulation.}

\begin{mathpar}
\inferrule*[right=T-if]{\vdash t : \tau \\ \vdash e : \tau}{b : \texttt{Bool} \vdash \texttt{if}(t, e)(b) : \tau}
\end{mathpar}

The iterator is operationally characterized by $\beta$ and $\eta$ laws. The $\beta$ laws say that
\texttt{if} should ``case correctly'': it should evaluate to its ``then'' branch on a \texttt{true}
input, and to its ``else'' branch on a \texttt{false} one. The $\eta$ law, meanwhile, says that
``things that behave like \texttt{if}s are judgmentally equal to \texttt{if}s;'' one can think of
this as a restricted form of function extensionality:
% TODO: can one so think of it?

\begin{mathpar}
% TODO: contexts? typing judgments?
\inferrule*[right=$\beta^{\texttt{if}}_1$]
  {\vdash t : \tau \\ \vdash e : \tau}{\vdash \texttt{if}(t, e)(\texttt{true}) \equiv t : \tau}

\inferrule*[right=$\beta^{\texttt{if}}_2$]
  {\vdash t : \tau \\ \vdash e : \tau}{\vdash \texttt{if}(t, e)(\texttt{false}) \equiv e : \tau}

\inferrule*[right=$\eta^{\texttt{if}}$]
  {x : \texttt{Bool} \vdash d : \tau \quad
   \vdash d[\texttt{true}/x] \equiv t : \tau \quad
   \vdash d[\texttt{false}/x] \equiv e : \tau}
  {\vdash d[b/x] \equiv \texttt{if}(t,e)(b) : \tau}
\end{mathpar}

As in our previous translations of equational laws, the $\beta$ laws correspond to equalities of
morphisms, while the $\eta$ law is a uniqueness criterion. In particular, the $\beta$ laws require
that if we have morphisms $t, e : 1 \to \llbr{\tau}$, then $\mathtt{if}(t,e) \circ \mathtt{true} =
t$ and $\mathtt{if}(t,e) \circ \mathtt{false} = e$. The interpretation of the $\eta$ law is more
involved, but it amounts to the requirement that $\mathtt{if}(t,e)$ be the \emph{unique} morphism
satisfying these properties. Putting these together, we obtain the following pair of commutative
triangles:

% https://q.uiver.app/#q=WzAsNCxbMCwwLCIxIl0sWzIsMCwiXFx0ZXh0dHR7Qm9vbH0iXSxbNCwwLCIxIl0sWzIsMiwiXFx0YXUiXSxbMCwxLCJcXHRleHR0dHt0cnVlfSJdLFsyLDEsIlxcdGV4dHR0e2ZhbHNlfSIsMl0sWzEsMywiXFx0ZXh0dHR7aWZ9KHQsZSkiLDJdLFswLDMsInQiLDFdLFsyLDMsImUiLDFdXQ==
\[\begin{tikzcd}
	1 && {\llbr{\texttt{Bool}}} && 1 \\
	\\
	&& \llbr{\tau}
	\arrow["{\texttt{true}}", from=1-1, to=1-3]
	\arrow["t"{description}, from=1-1, to=3-3]
	\arrow["{\texttt{if}(t,e)}"', from=1-3, to=3-3, dashed]
	\arrow["{\texttt{false}}"', from=1-5, to=1-3]
	\arrow["e"{description}, from=1-5, to=3-3]
\end{tikzcd}\]

For reasons that will become clear shortly, it is helpful to expand these triangles into squares:

% https://q.uiver.app/#q=WzAsNixbMCwwLCIxIl0sWzIsMCwiXFx0ZXh0dHR7Qm9vbH0iXSxbNCwwLCIxIl0sWzIsMiwiXFx0YXUiXSxbMCwyLCIxIl0sWzQsMiwiMSJdLFswLDEsIlxcdGV4dHR0e3RydWV9Il0sWzIsMSwiXFx0ZXh0dHR7ZmFsc2V9IiwyXSxbMSwzLCJcXHRleHR0dHtpZn0odCxlKSIsMl0sWzQsMywidCIsMV0sWzUsMywiZSIsMV0sWzIsNSwiXFxtYXRoc2Z7aWR9XzEiLDFdLFswLDQsIlxcbWF0aHNme2lkfV8xIiwxXV0=
\begin{equation}
\begin{tikzcd}
	1 && {\llbr{\texttt{Bool}}} && 1 \\
	\\
	1 && \llbr{\tau} && 1
	\arrow["{\texttt{true}}", from=1-1, to=1-3]
	\arrow["{\mathsf{id}_1}"{description}, from=1-1, to=3-1]
	\arrow["{\texttt{if}(t,e)}"', from=1-3, to=3-3, dashed]
	\arrow["{\texttt{false}}"', from=1-5, to=1-3]
	\arrow["{\mathsf{id}_1}"{description}, from=1-5, to=3-5]
	\arrow["t"{description}, from=3-1, to=3-3]
	\arrow["e"{description}, from=3-5, to=3-3]
\end{tikzcd}
\label{fig:bool-squares}
\end{equation}

\begin{exercise}
\sloppy % TODO: better workaround? (cf https://tex.stackexchange.com/q/153872)
Verify that $1 \oplus 1$ is a suitable choice for $\llbr{\texttt{Bool}}$. What are the denotations
of \texttt{true}, \texttt{false}, and $\texttt{if}(t,e)$?
\end{exercise}

\subsection{Categorifying Natural Numbers}

We now perform a similar translation for natural numbers. The type \texttt{Nat} has two
constructors, which we abbreviate as \texttt{z} (zero) and \texttt{s} (successor), and is equipped
with an iterator $\texttt{iter}(e_z, x.~e_s)(n)$. Notice that, unlike \texttt{if}, the iterator for
\texttt{Nat} binds a variable in the successor case: this will be the recursive result of evaluating
the iterator on the substructure of $n$.\footnote{We assume the reader is familiar with the standard
setup of recursive types. See, e.g., Chapter 15 of Harper, \emph{Practical Foundations for
Programming Languages (2\textsuperscript{nd} ed.)}.} The typing rules are given below:

\begin{mathpar}
\inferrule*[right=T-z]{~}{\vdash z : \mathtt{Nat}}

\inferrule*[right=T-s]{~}{n : \mathtt{Nat} \vdash s(n) : \mathtt{Nat}}

\inferrule*[right=T-iter]
  {\vdash e_z : \tau \\ x : \texttt{Nat} \vdash e_s : \tau}
  {n : \texttt{Nat} \vdash \mathtt{iter}(e_z, x. e_s)(n) : \tau}
\end{mathpar}

The $\beta$ rules for the iterator express its behavior as a fold over the structure of a natural
number: applying the iterator to zero evaluates to the first arm, while applying it to a successor
binds $x$ to the recursive result of folding over its predecessor, then evaluates $e_s$. As before,
we will elide detailed analysis of the $\eta$ rule beyond noting that it corresponds to a uniqueness
condition:

\begin{mathpar}
\inferrule*[right=$\beta^{\texttt{iter}}_z$]
  {\vdash e_z : \tau \\ x : \texttt{Nat} \vdash e_s : \tau}
  {\vdash \mathtt{iter}(e_z, x. e_s)(z) \equiv e_z : \tau}

\inferrule*[right=$\beta^{\texttt{iter}}_s$]
  {\vdash e_z : \tau \\ x : \texttt{Nat} \vdash e_s : \tau \\ \vdash n : \texttt{Nat}}
  {\vdash \mathtt{iter}(e_z, x. e_s)(\mathtt{s}(n)) \equiv e_s[\mathtt{iter}(e_z, x. e_s)(n)/x] : \tau}
\end{mathpar}

% TODO: a bit too metaphorical/scare-quote-y
As with Booleans, categorifying the rules for natural numbers yields an object representing the type
\texttt{Nat}, a morphism for each constructor, and, for each pair of ``folding-function'' morphisms
running ``parallel to'' the constructors, a unique morphism denoted by the iterator. In the case of
natural numbers, however, the resulting construction has a special name: the \emph{natural numbers
object} in a Cartesian closed category.\footnote{There are technical reasons why we insist that the
category be Cartesian closed, even though the construction does not require exponential objects.}

\begin{definition}[Natural numbers object]
\sloppy % TODO: better workaround? (cf https://tex.stackexchange.com/q/153872)
  A \emph{natural numbers object} in a Cartesian closed category $\calC$ consists of
  \begin{itemize}
    \item An object $N$ of $\calC$ and
    \item Morphisms $z : 1 \to N$ and $s : N \to N$
  \end{itemize}
  such that, for every object $A$ of $\calC$ and pair of morphisms $f_z : 1 \to A$ and $f_s : A \to
  A$, there exists a unique morphism $u$ such that the following diagram commutes:
  % https://q.uiver.app/#q=WzAsNSxbMCwwLCIxIl0sWzIsMCwiTiJdLFs0LDAsIk4iXSxbMiwyLCJBIl0sWzQsMiwiQSJdLFswLDEsInoiXSxbMSwzLCJ1IiwyXSxbNCwzLCJlX3MiLDFdLFsyLDQsInUiLDFdLFsyLDEsInMiLDJdLFswLDMsImVfeiIsMV1d
  \[\begin{tikzcd}
    1 && N && N \\
    \\
    && A && A
    \arrow["z", from=1-1, to=1-3]
    \arrow["{e_z}"{description}, from=1-1, to=3-3]
    \arrow["u"', from=1-3, to=3-3, dashed]
    \arrow["s"', from=1-5, to=1-3]
    \arrow["u"{description}, from=1-5, to=3-5, dashed]
    \arrow["{e_s}"{description}, from=3-5, to=3-3]
  \end{tikzcd}\]
\end{definition}

As before, it will be helpful to instead view this diagram as a pair of squares:

% https://q.uiver.app/#q=WzAsNyxbMCwwLCIxIl0sWzIsMCwiXFxsbGJye1xcdGV4dHR0e05hdH19Il0sWzQsMCwiXFxsbGJye1xcdGV4dHR0e05hdH19Il0sWzIsMiwiXFxsbGJye1xcdGF1fSJdLFswLDIsIjEiXSxbNCwyLCJcXGxsYnJ7XFx0YXV9Il0sWzMsMF0sWzAsMSwiXFx0ZXh0dHR7en0iXSxbMSwzLCJcXHRleHR0dHtyZWN9KGVfeix4Ln5lX3MpIiwyXSxbNCwzLCJlX3oiLDFdLFs1LDMsImVfcyIsMV0sWzIsNSwiXFx0ZXh0dHR7cmVjfShlX3oseC5+ZV9zKSIsMV0sWzAsNCwiXFxtYXRoc2Z7aWR9XzEiLDFdLFsyLDEsIlxcdGV4dHR0e3N9IiwyXV0=
\begin{equation}
\begin{tikzcd}
	1 && {\llbr{\texttt{Nat}}} & {} & {\llbr{\texttt{Nat}}} \\
	\\
	1 && {\llbr{\tau}} && {\llbr{\tau}}
	\arrow["{\texttt{z}}", from=1-1, to=1-3]
	\arrow["{\mathsf{id}_1}"{description}, from=1-1, to=3-1]
	\arrow["{\texttt{iter}(e_z,x.~e_s)}"', from=1-3, to=3-3, dashed]
	\arrow["{\texttt{s}}"', from=1-5, to=1-3]
	\arrow["{\texttt{iter}(e_z,x.~e_s)}"{description}, from=1-5, to=3-5, dashed]
	\arrow["{e_z}"{description}, from=3-1, to=3-3]
	\arrow["{e_s}"{description}, from=3-5, to=3-3]
\end{tikzcd}
\label{fig:nat-squares}
\end{equation}

\begin{exercise}
\sloppy % TODO: better workaround? (cf https://tex.stackexchange.com/q/153872)
  Give a categorical interpretation of the type of \emph{extended natural numbers}, i.e., the type
  of natural numbers extended with an extra element representing infinity, defined as follows:
\begin{lstlisting}
inductive ExtNat : Type where
  | zero : ExtNat
  | succ : ExtNat -> ExtNat
  | inf  : ExtNat
\end{lstlisting}

  \emph{Hint:} The resulting diagram will look a bit different from the two preceding examples.
\label{ex:ext-nat-interp}
\end{exercise}

\subsection{A General Categorification of Inductive Types}

Let us now take stock of the preceding examples and generalize what we have seen. In each of our
translations above, we did the following:

\begin{itemize}
  \item First, interpret the constructors of an inductive type as morphisms from the interpretation
  of the constructor's domain to the type's denotation. That is, a constructor $\texttt{c} :
  \texttt{D} \to \texttt{T}$ of \texttt{T} becomes a morphism $\llbr{\texttt{D}}
  \xrightarrow{\texttt{c}} \llbr{\texttt{T}}$.
  \item Then, interpret the iterator as a unique morphism $\llbr{\texttt{T}} \to A$ for every
  collection of morphisms into $A$ ``parallel to the constructors,'' i.e., morphisms $t_i$ that map
  into $A$ from the domain of the $i$th constructor where $\llbr{\texttt{T}}$ has been replaced by
  $A$. For instance, if $\texttt{T}$ has two constructors
  \[ \texttt{c}_1 : \texttt{Bool} \to \texttt{T} \quad\text{and}\quad
     \texttt{c}_2 : \texttt{T} \times \texttt{Nat} \to \texttt{T}, \]
  its iterator is interpreted as a unique morphism $u_{f_1, f_2} : \llbr{\texttt{T}} \to A$ for
  all objects $A$ and morphisms
  \[ f_1 : \llbr{\texttt{Bool}} \to A \quad\text{and}\quad
     f_2 : A \times \llbr{\texttt{Nat}} \to A. \]
  Moreover, the resulting squares must commute.
\end{itemize}

This procedure is a useful heuristic but is not yet a precise categorical account of an inductive
type. For instance, the ``replacement'' performed in the second step is merely a syntactic
operation. We will now make this categorical translation more precise.

We start by consolidating our commutative squares. Recall that any inductive type with multiple
constructors can be equivalently expressed as a having a single constructor whose domain is the sum
of the domains of the original constructors. That is, the constructors $c_1 : A \to T$, $c_2 : B \to
T$, and $c_3 : C \to T$ can be equivalently represented by a single constructor $c : A \oplus B
\oplus C \to T$. The iterator is then parameterized by a single morphism running  ``parallel'' to
the constructor that maps into $A$ from an analogous coproduct over $A$.

Applying this rewriting to the \texttt{Nat} datatype depicted in diagram
(\ref{fig:nat-squares}), we obtain
% https://q.uiver.app/#q=WzAsNCxbMCwwLCIxIFxcb3BsdXMgXFxsbGJye1xcdGV4dHR0e05hdH19Il0sWzIsMCwiXFxsbGJye1xcdGV4dHR0e05hdH19Il0sWzAsMiwiMSBcXG9wbHVzIEEiXSxbMiwyLCJBIl0sWzAsMSwiXFxsYW5nbGUgeiwgc1xccmFuZ2xlIl0sWzEsMywiXFxtYXRodHR7aXRlcn0oXFxsYW5nbGUgZl96LCBmX3NcXHJhbmdsZSkiXSxbMiwzLCJcXGxhbmdsZSBmX3osIGZfc1xccmFuZ2xlIiwyXSxbMCwyLCJcXGxlZnRcXGxhbmdsZSBcXG1hdGhzZntpbmx9IFxcY2lyYyBcXG1hdGhzZntpZH1fMSwgXFxtYXRoc2Z7aW5yfSBcXGNpcmMgXFxtYXRoc2Z7aXRlcn0oXFxsYW5nbGUgZl96LCBmX3NcXHJhbmdsZSkgXFxyaWdodFxccmFuZ2xlIiwyXV0=
\[\begin{tikzcd}
	{1 \oplus \llbr{\texttt{Nat}}} && {\llbr{\texttt{Nat}}} \\
	\\
	{1 \oplus A} && A
	\arrow["{\langle z, s\rangle}", from=1-1, to=1-3]
	\arrow["{\left\langle \mathsf{inl} \mathop\circ \mathsf{id}_1,~ \mathsf{inr} \mathop\circ \mathsf{iter}(\langle f_z, f_s\rangle) \right\rangle}"', from=1-1, to=3-1]
	\arrow["{\mathtt{iter}(\langle f_z, f_s\rangle)}", from=1-3, to=3-3]
	\arrow["{\langle f_z, f_s\rangle}"', from=3-1, to=3-3]
\end{tikzcd}\]
and a similar transformation occurs for diagram (\ref{fig:bool-squares}). (In the diagram above, we
are using the notation $\langle f, g\rangle$ to denote the unique morphism obtained from the
universal property of coproducts. Thinking of morphisms as functions, $\langle f, g\rangle$ does a
case split on its input, mapping elements of the left-hand summand using $f$ and those of the
right-hand summand using $g$.)

The top-left corner should look familiar: it is the body of the recursive type equation that defines
the natural numbers. That is, using a least-fixed-point operator $\mu$, we would write
\[ \texttt{Nat} \triangleq \mu t.~1 \oplus t, \]
so that \texttt{Nat} is the (smallest) type satisfying the type isomorphism
\[ \texttt{Nat} \cong 1 \oplus \texttt{Nat}. \]

\begin{exercise}
\sloppy % TODO: better workaround? (cf https://tex.stackexchange.com/q/153872)
  ``Consolidate'' diagram (\ref{fig:bool-squares}) as we did for diagram (\ref{fig:nat-squares})
  above, and verify that the (degenerate) recursive type equation for Booleans emerges.
\end{exercise}

With this consolidation, every inductive type now translates to a single object (the top-right
corner of the square), a single morphism (the constructor), and a universal property that translates
the iterator. The ``shape'' of the objects on the left-hand side of the commutative square are
determined by the ``shape'' of the inductive type. More precisely, using the observation in the
preceding paragraph, that ``shape'' is determined precisely by the expression $\tau$ in the
fixed-point equation $\mu t.~\tau$ for the inductive type.

Adopting a categorical lens, we might observe that an expression $t.~\tau$ defines a mapping from a
type $t$ to some type $\tau$ parameterized thereby. In our categorical interpretation, types form a
category; thus, such a mapping from types to types ought to correspond to a functor. More
specifically, since we are mapping from a category into itself, this is an \emph{endofunctor}. The
endofunctor associated with the type of natural numbers is defined by
\begin{align*}
  F_\texttt{Nat}(X) &= 1 \oplus X\\
  F_\texttt{Nat}(f : X \to Y) &= \langle \mathsf{inl} \circ \mathsf{id}_1, \mathsf{inr} \circ f\rangle : 1 \oplus X \to 1 \oplus Y
\end{align*}

Of course, the mapping $t.~\tau$ is not itself a type; rather, we apply the least-fixed-point
operator $\mu$ to obtain an inductive type. We should expect a corresponding maneuver in the world
of categories: we interpret an inductive type as the \emph{least fixed point} $\mu F$ of an
endofunctor $F$ on our semantic domain: i.e., the ``least'' object $\mu F$ satisfying $F(\mu F) =
\mu F$.

This naturally invites the following question: in the absence of an \emph{a priori} ordering on
objects, how do we make precise the fact that $\mu F$ should be ``least'' among the fixed points of
$F$? Appealing to our categorical intuition, we should expect this to be some kind of initiality or
colimiting property (recall that, using our order-theoretic lens, we think of initial objects as
``minimal''). And, in fact, as we will now demonstrate, we have already seen such a property in our
Boolean and natural-number examples!

The key is to recall that an inductive type is defined not only by an object but also by a morphism
representing its constructor that tells us how to ``map into'' the type. Thus, instead of merely
considering the ``smallest'' object among the fixed points of $F$, we should instead consider the
``smallest'' \emph{pair} of an $F$-fixed point \emph{and} its corresponding constructor morphism.
Such a pair is an example of an \emph{algebra} over the endofunctor $F$, also known as an
\emph{$F$-algbera}.

\begin{definition}[$F$-Algebra]
\sloppy % TODO: better workaround? (cf https://tex.stackexchange.com/q/153872)
  Let $F : \calC \to \calC$ be an endofunctor. An \emph{$F$-algebra}, or \emph{algebra over $F$}, is
  a pair $(A, \alpha)$ where $A$ is an object in $\calC$ and $\alpha : F A \to A$ is a morphism.
  
  $A$ is called the \emph{carrier} of the $F$-algebra, and $\alpha$ is the \emph{structure map}.
  Note that the structure map $\alpha$ alone is sometimes referred to as an $F$-algebra, leaving $A$
  implicit.
\end{definition}

We are looking for the ``smallest'' $F$-algebra $(A, \alpha)$ such that its carrier $A$ is a fixed
point. Let us first, however, consider simply the ``smallest'' $F$-algebra, without any further
restrictions. We do this because it turns out that there is a natural interpretation of ``smallest''
in terms of initiality. In particular, $F$-algebras form a category $\mathsf{Alg}_F$, so the
``smallest'' $F$-algebra is simply the initial object in this category.

\begin{definition}[Category of $F$-algebras]
  Fix an endofunctor $F : \calC \to \calC$. The \emph{category of $F$-algebras} (or \emph{category
  of algebras over $F$}) is defined by the following:
  \begin{itemize}
    \item Objects are $F$-algebras $(A, \alpha : F A \to A)$.
    \item Morphisms from $(A, \alpha)$ to $(B, \beta)$ are morphisms $f : A \to B$ in $\calC$ such
          that the following square commutes:
    % https://q.uiver.app/#q=WzAsNCxbMCwwLCJGIEEiXSxbMiwwLCJBIl0sWzAsMiwiRiBCIl0sWzIsMiwiQiJdLFswLDEsIlxcYWxwaGEiXSxbMSwzLCJmIl0sWzIsMywiXFxiZXRhIiwyXSxbMCwyLCJGIGYiLDJdXQ==
    \[\begin{tikzcd}
      {F A} && A \\
      \\
      {F B} && B
      \arrow["\alpha", from=1-1, to=1-3]
      \arrow["{F f}"', from=1-1, to=3-1]
      \arrow["f", from=1-3, to=3-3]
      \arrow["\beta"', from=3-1, to=3-3]
    \end{tikzcd}\]
    \item Identity and composition are lifted from $\calC$.
  \end{itemize}
\end{definition}

\begin{definition}[Initial $F$-algebra]
\sloppy % TODO: better workaround? (cf https://tex.stackexchange.com/q/153872)
  Fix an endofunctor $F : \calC \to \calC$. The \emph{initial $F$-algebra} (or \emph{initial algebra
  over $F$}), if it exists, is the initial object in the category of $F$-algebras.

  We write $\mathsf{in}_F$ for the (structure map of) the initial $F$-algebra. Given another
  $F$-algebra $(A, \alpha)$, we write $\llpr{\alpha}_F$ (or just $\llpr{\alpha}$ if $F$ is clear
  from context) for the unique morphism from the initial object to $(A, \alpha)$.
\end{definition}

We have thus found a suitable notion of a ``smallest'' $F$-algebra. But what about the smallest
$F$-algebra whose carrier is a fixed point of $F$? In fact, these are one and the same (up to
isomorphism): it will always be the case that the carrier of the initial $F$-algebra is (isomorphic
to) a fixed point of $F$. This is the content of \emph{Lambek's lemma}.

\begin{proposition}[Lambek's lemma]
  Let $F : \calC \to \calC$ be an endofunctor. If $(A, \mathsf{in}_F : F A \to A)$ is the initial
  $F$-algebra, then $\mathsf{in}_F$ is an isomorphism.
\end{proposition}

\begin{proof}
  Let $(A, \mathsf{in}_F : F A \to A)$ be the initial $F$-algebra. We will show that $\mathsf{in}_F$
  is an isomorphism by showing that its inverse is $\llpr{F~\mathsf{in}_F} : A \to F A$, the unique
  $F$-algebra morphism $(A, \mathsf{in}_F) \to (F A, F~\mathsf{in}_F)$ induced by initiality.

  First, observe that $\mathsf{in}_F$ is an $F$-algebra morphism $(F A, F~\mathsf{in}_F) \to (A,
  \mathsf{in}_F)$, since the following diagram trivially commutes:

  % https://q.uiver.app/#q=WzAsNCxbMCwwLCJGIChGIEEpIl0sWzIsMCwiRiBBIl0sWzAsMiwiRiBBIl0sWzIsMiwiQSJdLFswLDEsIkYgXFxhbHBoYSJdLFswLDIsIkYgXFxhbHBoYSIsMl0sWzEsMywiXFxhbHBoYSJdLFsyLDMsIlxcYWxwaGEiLDJdXQ==
  \[\begin{tikzcd}
    {F (F A)} && {F A} \\
    \\
    {F A} && A
    \arrow["{F~\mathsf{in}_F}", from=1-1, to=1-3]
    \arrow["{F~\mathsf{in}_F}"', from=1-1, to=3-1]
    \arrow["\mathsf{in}_F", from=1-3, to=3-3]
    \arrow["\mathsf{in}_F"', from=3-1, to=3-3]
  \end{tikzcd}\]

  Thus, we have the following diagram in $\mathsf{Alg}_F$:

  % https://q.uiver.app/#q=WzAsMyxbMCwwLCIoQSwgXFxhbHBoYSkiXSxbMCwyLCIoRiBBLCBGIFxcYWxwaGEpIl0sWzAsNCwiKEEsIFxcYWxwaGEpIl0sWzAsMSwiXFxsbHBye0YgXFxhbHBoYX0iXSxbMSwyLCJcXGFscGhhIl0sWzAsMiwiXFxtYXRoc2Z7aWR9X3soQSwgXFxhbHBoYSl9IiwyLHsiY3VydmUiOjR9XV0=
  \[\begin{tikzcd}
    {(A, \mathsf{in}_F)} \\
    \\
    {(F A, F~\mathsf{in}_F)} \\
    \\
    {(A, \mathsf{in}_F)}
    \arrow["{\llpr{F~\mathsf{in}_F}}", from=1-1, to=3-1]
    \arrow["{\mathsf{id}_A}"', curve={height=36pt}, from=1-1, to=5-1]
    \arrow["\mathsf{in}_F", from=3-1, to=5-1]
  \end{tikzcd}\]
  
  Since $(A, \mathsf{in}_F)$ is initial, there is only one $F$-algebra morphism from $(A,
  \mathsf{in}_F)$ to itself, so this diagram must commute, i.e.,
  \[ \mathsf{in}_F \circ \llpr{F~\mathsf{in}_F} = \mathsf{id}_A.\]
  Accordingly, $\llpr{F~\mathsf{in}_F}$ is a right-hand inverse for $\mathsf{in}_F$.
  
  To show that it is also a left-hand inverse, we apply $F$ to both sides of the preceding equation
  and rewrite using the functor laws:
  \[ F~\mathsf{in}_F \circ F \llpr{F~\mathsf{in}_F} = \mathsf{id}_{F A}. \]
  Then we observe that
  \[ F~\mathsf{in}_F \circ F \llpr{F~\mathsf{in}_F} = \mathsf{in}_F \circ \llpr{F~\mathsf{in}_F} \]
  from the fact that $\llpr{F~\mathsf{in}_F}$ is an $F$-algebra morphism:
  % https://q.uiver.app/#q=WzAsNCxbMCwwLCJGIEEiXSxbMiwwLCJBIl0sWzAsMiwiRiAoRiBBKSJdLFsyLDIsIkYgQSJdLFswLDEsIlxcbWF0aHNme2lufV9GIl0sWzIsMywiRn5cXG1hdGhzZntpbn1fRiJdLFswLDIsIkYgXFxsbHBye0Z+XFxtYXRoc2Z7aW59X0Z9IiwxXSxbMSwzLCJcXGxscHJ7Rn5cXG1hdGhzZntpbn1fRn0iLDFdXQ==
  \[\begin{tikzcd}
    {F A} && A \\
    \\
    {F (F A)} && {F A}
    \arrow["{\mathsf{in}_F}", from=1-1, to=1-3]
    \arrow["{F \llpr{F~\mathsf{in}_F}}"{description}, from=1-1, to=3-1]
    \arrow["{\llpr{F~\mathsf{in}_F}}"{description}, from=1-3, to=3-3]
    \arrow["{F~\mathsf{in}_F}", from=3-1, to=3-3]
  \end{tikzcd}\]
  So, substituting the second equation into the first, we have
  \[ \mathsf{in}_F \circ \llpr{F~\mathsf{in}_F} = \mathsf{id}_{F A}, \]
  completing the proof that $\mathsf{in}_F : F A \to A$ is an isomorphism.
\end{proof}

With Lambek's lemma in hand, we can now give a precise categorification of inductive types: the
inductive type represented by $\mu t.~\tau$ is interpreted as the initial algebra of the endofunctor
$F(t) = \llbr{\tau}$, where we extend our previous translation of types to map the free type
variable $t$ to the functor variable of the same name.\footnote{The action of $F$ on morphisms must
also be defined in the ``obvious'' way. We elide the details, but one can inspect the definition of
$F_\texttt{Nat}$ at the beginning of this subsection for an example.} The algebra's carrier $\mu F$
is the object representing the type, while its structure map $\mathsf{in}_F$ represents the type's
constructor.

It remains only to precisify our former account of iterators. Recall our previous observation that
iterators map out of an inductive type $\mu F$ into some other type $A$ using a ``combining
morphism'' that runs ``parallel'' to the constructor morphism---i.e., one that maps into $A$ from an
object that ``replaces'' instances of $\mu F$ with instances of $A$. Using our newfound language of
functors, we can phrase this much more concisely: an iterator maps from $\mu F$ to some type $A$
using a ``combining morphism'' $f_c : F A \to A$. Or, as a diagram:

% https://q.uiver.app/#q=WzAsNCxbMCwwLCJGKFxcbXUgRikiXSxbMiwwLCJcXG11IEYiXSxbMCwyLCJGIEEiXSxbMiwyLCJBIl0sWzAsMSwiYyJdLFsyLDMsImZfYyJdLFswLDIsIkYgKFxcbWF0aHNme2l0ZXJ9X3tmX2N9KSIsMV0sWzEsMywiXFxtYXRoc2Z7aXRlcn1fe2ZfY30iLDFdXQ==
\[\begin{tikzcd}
	{F(\mu F)} && {\mu F} \\
	\\
	{F A} && A
	\arrow["c", from=1-1, to=1-3]
	\arrow["{F (\mathsf{iter}_{f_c})}"{description}, from=1-1, to=3-1]
	\arrow["{\mathsf{iter}_{f_c}}"{description}, from=1-3, to=3-3]
	\arrow["{f_c}", from=3-1, to=3-3]
\end{tikzcd}\]

In fact, we already saw this diagram (diagrams (\ref{fig:bool-squares}) and (\ref{fig:nat-squares})
are variations of it in the particular cases of natural numbers and Booleans). But having now
developed the vocabulary of $F$-algebras, this should look familiar for another reason: iterators
are $F$-algebra morphisms! More precisely, given an inductive type $\mu F$, the ``unapplied''
iterator function is represented by $\llpr{-}_F$; given a particular ``combining morphism'' $f_c : F
A \to A$---which is precisely an $F$-algebra over $A$---the resulting parameterized iterator is
$\llpr{f_c} : \mu F \to A$.

As an example, we translate the parity function on natural numbers, which returns \texttt{true} if
its input is even and \texttt{false} if it is odd, as an application of the natural-number iterator.
First, we identify the appopriate ``combining function'' for this fold. Since we are ultimately
producing a Boolean, this function should translate to a morphism of type $F_\texttt{Nat}
\llbr{\texttt{Bool}} \to \llbr{\texttt{Bool}}$. We first define its programmatic version
below:\footnote{We assume predefined operations on Booleans}

\begin{align*}
\texttt{f}_\texttt{even} &: 1 \oplus \texttt{Bool} \to \texttt{Bool}\\
\texttt{f}\texttt{even}~(\mathsf{inl}~()) &= \texttt{true}\\
\texttt{f}_\texttt{even}~(\mathsf{inr}~\texttt{b}) &= \texttt{not}(\texttt{b})
\end{align*}

This becomes the $F_\texttt{Nat}$-algebra
% TODO: does this type-check?
\[ \alpha_\texttt{even} = \langle \mathtt{true}, \mathtt{not} \circ - \rangle \] which finally
yields our translation of the parity function $\llpr{\alpha_\texttt{even}} : \llbr{\texttt{Nat}} \to
\llbr{\texttt{Bool}}$.

\begin{exercise}
\sloppy % TODO: better workaround? (cf https://tex.stackexchange.com/q/153872)
  Fill in the missing piece of the above translation by categorifying \texttt{not} using a suitable
  algebra over the endofunctor characterizing Booleans.
\end{exercise}

To put together the developments in this section, let's consider one more example of a categorified
inductive datatype, now expressed in terms of initial algebras. Recall the inductive datatype of
lists:

\begin{lstlisting}
inductive List (A : Type) : Type where
  | nil : List A
  | cons : A -> List A -> List A
\end{lstlisting}

Notice that this type is parameterized over a fixed type \texttt{A}; we will write $A$ for that
type's categorical denotation. We can read the endofunctor of which \texttt{List A} is a least fixed
point off of its constructors:\footnote{We again elide the action on morphisms.}
\[ F_\mathtt{List_A}(X) = 1 \oplus A \times X. \]
Lists are thus represented by the initial algebra $(\mu F_{\mathtt{List_A}},
\mathsf{in}_{F_\mathtt{List_A}})$. Returning to our example from earlier, we can define the list
length function as the fold induced by the ``combining function''
\begin{align*}
  \texttt{f}_\texttt{length} &: \texttt{Unit} \oplus \texttt{A} \times \texttt{Nat} \to \texttt{Nat}\\
  \texttt{f}_\texttt{length}~(\mathsf{inl}~()) &= 0\\
  \texttt{f}_\texttt{length}~(\mathsf{inr}~(a, \ell)) &= \texttt{succ}(\ell)
\end{align*}
which we translate into the $F_\mathtt{List_A}$-algebra
% TODO: does this type-check?
\[ \alpha_\texttt{length} = \langle z, s \circ \pi_1 \rangle \]
to obtain the fold $\llpr{\alpha_\texttt{length}} : \mu F_\mathtt{List_A} \to \llbr{\texttt{Nat}}$.

\begin{exercise}
Categorify the type of \emph{rose trees} defined below:
\begin{lstlisting}
inductive Rose (A : Type) : Type where
  | node : A -> List Rose -> Rose
\end{lstlisting}
TODO: finish
\end{exercise}

Up to this point, we have simply assumed that, given an endofunctor $F$, we can always find an
initial $F$-algebra. This is not, however, always the case. And we should expect as much: after all,
not all recursive type equations admit a least fixed point---this is why we insist that no
non-positive occurrences of $t$ appear in $\tau$ when defining an inductive type $\mu t.~\tau$. The
failure to obtain a least fixed point in a recursive type equation corresponds to an analogous lack
of an initial algebra for the corresponding functor. For instance, $\mu t.~t \to 2$ does not admit a
least fixed point, motivating the following exercise:

\begin{exercise}
Show that the endofunctor $F(X) = 2^X$ on $\Set$ does not admit an initial algebra.
(\emph{Hint:} recall Lambek's lemma.)
\end{exercise}

Of course, such results generalize to arbitrary categories, not just $\Set$, but examining
them in depth would take us too far afield. It suffices to note that all the usual strictly positive
inductive types of interest are characterized by endofunctors that admit initial algebras.

\begin{exercise}
\sloppy % TODO: better workaround? (cf https://tex.stackexchange.com/q/153872)

  Can we dualize this? (Spoiler: yes.) That is, given an endofunctor $F$, let us consider
  \emph{$F$-coalgebras} $(A, \eta : A \to F A)$ instead of $F$-algebras $(A, \alpha : F A \to A)$,
  and examine the \emph{terminal} $F$-coalgebra rather than the initial one.
  
  TODO: More scaffolding, give a specific functor example (e.g., streams)

  \begin{enumerate}
    \item What type does the carrier of the terminal $F$-coalgebra represent? (\emph{Hint}: It will
          turn out that the carrier of the terminal $F$-coalgebra is also a fixed point of $F$ up to
          isomorphism.) What about its structure map?

    \item How do we interpret the unique $F$-coalgebra morphism from another $F$-coalgebra into the
          terminal one?
  \end{enumerate}
\end{exercise}

\section{The Families Fibration: Indexing Inductive Types}

Having given a categorical translation of simple inductive types, we now turn our attention to
\emph{indexed} inductive types. Recall that indexed inductive types enrich inductive types by
allowing types to be parameterized by data values at the type level: for instance, in addition to
the simple type \texttt{List A} of arbitrary-length lists with elements of type \texttt{A}, we can
also define the indexed type \texttt{Vector A n} of \texttt{A}-containing lists of length
\texttt{n}.

% This is the only feature of a dependent type system we will require, so we give a categorical
% account specifically of indexed inductive types rather than full dependent types. The advantage of
% this simplification is that such an account follows naturally from our preceding work.

While we gave a generalized presentation of inductive types in terms of initial algebras of
endofunctors over arbitrary categories, hereafter we will specifically consider a categorical
semantics in terms of sets. While all of the ideas presented can be generalized to other categories,
focusing on $\Set$ will simplify some of the required categorical machinery. (Note that our
preceding interpretation of inductive types can indeed be carried out in $\Set$, since it is
Cartesian closed and contains finite coproducts---i.e., it is \emph{bicartesian closed}.)

The category $\Set$ alone, however, lacks the structure to give an elegant semantics to indexed
inductive types. Instead, we will turn to a new category built atop $\Set$ that is equipped with a
notion of ``indexing'' that will afford a straightforward interpretation of our inductive types'
indices.

\begin{definition}[Category of families (of sets)]
  The \emph{category of families (of sets)} $\FamSet$ is defined as follows:
  \begin{itemize}
    \item Objects are pairs $(A, P : A \to \Ob(\Set))$ of an indexing set $A$ and a \emph{family} of
          sets $P$.
    \item Morphisms from $(A, P)$ to $(B, Q)$ are pairs of functions $(f, f^\sim)$ such that $f : A
          \to B$ and $f^\sim : \prod_{a \in A} P(a) \to Q(f(a))$.
  \end{itemize}
\end{definition}

Just as inductive types were interpreted as initial algebras of endofunctors over $\Set$, indexed
inductive types can be interpreted as initial algebras of endofunctors over $\FamSet$. As an
example, consider the type \texttt{Fin n}, the type of ``natural numbers less than $n$'' (or,
isomorphically, $\Z/n\Z$):

\begin{lstlisting}
inductive Fin : Nat -> Type where
  | finZero : {n : Nat} -> Fin (succ n)
  | finSucc : {n : Nat} -> Fin n -> Fin (succ n)
\end{lstlisting}

The corresponding endofunctor on $\FamSet$

TODO:

\begin{exercise}
  Translate the type $\mathtt{Vector_A}$ TODO: finish
\end{exercise}

% TODO: I think we actually just use the families fibration
% It may not be a surprise that our categorical interpretation of \emph{indexed} inductive types will
% be in the category of indexed sets, which we now define.\footnote{This definition carries with it
% the usual provisos about Grothendieck universes, which we elide.}

% \begin{definition}[Category of indexed sets]
% \sloppy % TODO: better workaround? (cf https://tex.stackexchange.com/q/153872)
%   Given a set $I$, the category $\Set_I$ of $I$-indexed sets is defined as follows:
  
%   \begin{itemize}
%     \item Objects are $I$-indexed sets $\{S_i\}_{i \in I}$
%     \item Morphisms from $\{S_i\}_{i \in I}$ to $\{T_i\}_{i \in I}$ are set functions $f :
%     \prod_{i \in I} S_i \to T_i$
%     \item Composition and identities are defined using the corresponding set definitions at each
%     index
%     % TODO: should we spell out composition?
%     % \item Given $f : \{S_i\}_{i \in I} \to \{T_i\}_{i \in I}$ and $g : \{T_i\}_{i \in I} \to
%     % \{U_i\}_{i \in I}$, their composite $g \circ f : \{S_i\}_{i \in I} \to \{U_i\}_{i \in I}$ is
%     % defined by $i \mapsto g(i) \circ f(i)$
%     % \item Identities
%   \end{itemize}
% \end{definition}




\vspace{3cm}

\subsection{Indexing Inductive Types}

TODO: write exercise solutions!

TODO: texttt and mathttt are not the same --- settle on one

\section{Constructing Inductive Types from Refinements}

\section{Correctness}



\section{}

\section{}

References:
- Article
- Max New lecture notes (buggy)
- Wadler, Recursive types for free!
